% The Basel Problem — a slideSonnet demo deck.
%
% Compile:  slidesonnet sty           (drops slidesonnet.sty next to this file)
%           latexmk -pdf basel-problem.tex
% Narrate:  edit basel-problem.narration  (or: slidesonnet edit basel-problem.pdf)
% Render:   slidesonnet export basel-problem.pdf -o basel-problem.mp4 --engine piper
%
% Every emitted page carries a stable \ssid slide-id; the narration lives in
% basel-problem.narration, keyed by those ids.
\documentclass[aspectratio=169]{beamer}
\usepackage{slidesonnet}
\usetheme{Madrid}
\usecolortheme{seahorse}
\setbeamertemplate{navigation symbols}{}
\definecolor{accent}{HTML}{E94560}

\title{The Basel Problem}
\subtitle{An impossible sum, a young genius, and a surprise visit from $\pi$}
\author{A slideSonnet demo}
\date{}

\begin{document}

\begin{frame}[plain]
  \ssid{intro-title}
  \titlepage
\end{frame}

\begin{frame}{A simple question}
  \ssid{simple-question}
  What happens when you add the reciprocals of the perfect squares?
  \[ \frac{1}{1} + \frac{1}{4} + \frac{1}{9} + \frac{1}{16} + \frac{1}{25} + \cdots \]
  In modern notation:
  \[ \sum_{n=1}^{\infty} \frac{1}{n^2} = \ ? \]
\end{frame}

\begin{frame}{The challenge (1650)}
  \ssid<1>{challenge-mengoli}
  \ssid<2>{challenge-bernoulli-intro}
  \ssid<3>{challenge-bernoulli-quote}
  \textbf{Pietro Mengoli} posed this problem in Bologna in 1650.
  The sum converges to something near \textbf{1.6449}\ldots\ but \emph{what} is it?
  \onslide<2->{\medskip Even \textbf{Jakob Bernoulli} admitted defeat in 1689:}
  \onslide<3->{%
    \begin{quote}\color{accent}\itshape
      ``If anyone finds and communicates to us that which thus far has eluded
      our efforts, great will be our gratitude.''
    \end{quote}}
\end{frame}

\begin{frame}{Enter Leonhard Euler}
  \ssid<1>{enter-euler}
  \ssid<2>{euler-answer}
  The problem stayed open for \textbf{85 years}. Then, in 1734 in Saint
  Petersburg, a young mathematician from Basel named \textbf{Leonhard Euler}
  found the answer.
  \onslide<2->{%
    \[ \sum_{n=1}^{\infty} \frac{1}{n^2} = \frac{\pi^2}{6} \]
    $\pi$? In a sum that has nothing to do with circles?}
\end{frame}

\begin{frame}{Step 1: start with $\sin x$}
  \ssid<1>{step1-series}
  \ssid<2>{step1-divide}
  \ssid<3>{step1-coeff}
  The Taylor series for sine:
  \[ \sin x = x - \frac{x^3}{3!} + \frac{x^5}{5!} - \frac{x^7}{7!} + \cdots \]
  \onslide<2->{Divide both sides by $x$:
    \[ \frac{\sin x}{x} = 1 - \frac{x^2}{6} + \frac{x^4}{120} - \cdots \]}
  \onslide<3->{The coefficient of $x^2$ is $\color{accent}-\tfrac{1}{6}$ --- hold onto it.}
\end{frame}

\begin{frame}{Step 2: factor by the roots}
  \ssid<1>{step2-zeros}
  \ssid<2>{step2-factor}
  $\dfrac{\sin x}{x}$ has zeros at $x = \pm\pi, \pm 2\pi, \pm 3\pi, \ldots$
  \onslide<2->{Euler treated it like a polynomial and factored by its roots:
    \[ \frac{\sin x}{x} = \left(1 - \frac{x^2}{\pi^2}\right)
       \left(1 - \frac{x^2}{4\pi^2}\right)
       \left(1 - \frac{x^2}{9\pi^2}\right)\cdots \]
    \footnotesize (Justified much later by the Weierstrass factorization theorem, 1876.)}
\end{frame}

\begin{frame}{Step 3a: reading the $x^2$ coefficient}
  \ssid{step3a-finite}
  With $a_k = \frac{1}{k^2\pi^2}$, three factors give
  \[ (1 - a_1 x^2)(1 - a_2 x^2)(1 - a_3 x^2)
     = 1 - (a_1 + a_2 + a_3)\,x^2 + \cdots \]
  Each $x^2$ term picks $-a_k x^2$ from \emph{one} factor and $1$ from the rest.
\end{frame}

\begin{frame}{Step 3b: the infinite product}
  \ssid{step3b-infinite}
  \[ \prod_{n=1}^{\infty}\left(1 - \frac{x^2}{n^2\pi^2}\right)
     = 1 - \left(\frac{1}{\pi^2}\sum_{n=1}^{\infty}\frac{1}{n^2}\right)x^2 + \cdots \]
  So the $x^2$ coefficient is $-\dfrac{1}{\pi^2}\displaystyle\sum_{n=1}^{\infty}\frac{1}{n^2}$.
\end{frame}

\begin{frame}{Step 4: the punchline}
  \ssid<1>{punchline-taylor}
  \ssid<2>{punchline-product}
  \ssid<3>{punchline-equate}
  From the Taylor series: $\dfrac{\sin x}{x} = 1 \color{accent}- \tfrac{1}{6}\color{black}\,x^2 + \cdots$
  \onslide<2->{\medskip From the product:
    $\dfrac{\sin x}{x} = 1 \color{accent}- \tfrac{1}{\pi^2}\sum \tfrac{1}{n^2}\color{black}\,x^2 + \cdots$}
  \onslide<3->{\medskip Coefficients must agree:
    \[ \frac{1}{\pi^2}\sum_{n=1}^{\infty}\frac{1}{n^2} = \frac{1}{6} \]}
\end{frame}

\begin{frame}{The result}
  \ssid{final-result}
  Multiply both sides by $\pi^2$:
  \[ 1 + \frac{1}{4} + \frac{1}{9} + \frac{1}{16} + \cdots
     = \frac{\pi^2}{6} \approx 1.6449\ldots \]
  The circle constant was hiding in the sum all along.
\end{frame}

\begin{frame}{Why it matters}
  \ssid<1>{why-rigor}
  \ssid<2>{why-zeta}
  Euler's proof was not rigorous by modern standards --- but the answer was
  never in doubt.
  \onslide<2->{\medskip The sum is $\zeta(2)$, a value of the Riemann zeta function
    $\zeta(s) = \sum_{n=1}^{\infty} n^{-s}$. Even values are rational multiples of
    powers of $\pi$; the odd values remain mysterious.}
\end{frame}

\begin{frame}{What we learned}
  \ssid{what-learned}
  \[ 1 + \frac{1}{4} + \frac{1}{9} + \frac{1}{16} + \cdots = \frac{\pi^2}{6} \]
  \begin{itemize}
    \item A simple question can hide a profound answer.
    \item The boldness to treat $\sin x / x$ as a polynomial was the key.
    \item $\pi$ connects geometry and number theory in ways we still explore.
  \end{itemize}
\end{frame}

\begin{frame}[plain]
  \ssid{sources}
  \centering \large Thanks for watching.
\end{frame}

\end{document}
